% minimal.tex — the smallest valid rrxiv paper.
%
% One claim, one piece of evidence, one citation, one open question.
% Used as a parser conformance fixture: the rrxiv-python parser must
% produce a CIR from this file that validates against cir.schema.json.
%
% Compile:
%   tectonic --keep-intermediates minimal.tex
%   # or: pdflatex minimal.tex && pdflatex minimal.tex

\documentclass{rrxiv}

\rrxivid{rrxiv-example-minimal}
\rrxivversion{v1}
\rrxivprotocolversion{0.1.0}
\rrxivlicense{CC-BY-4.0}
\rrxivtopics{example,conformance}

\title{A minimal rrxiv paper}
\author{rrxiv Project}
\date{2026-05-04}

\begin{document}
\maketitle

\begin{abstract}
This paper exists as a parser conformance fixture. It contains one claim,
one piece of evidence, one citation, and one open question. Its only
purpose is to round-trip through \texttt{rrxiv-python}'s parser into a
valid Canonical Intermediate Representation (CIR).
\end{abstract}

\section{The claim}
\label{sec:claim}

\begin{claim}[Conformance fixture]
\label{claim:fixture}
A minimal rrxiv paper, when compiled by \texttt{rrxiv.cls} v0.1, emits a
sidecar \texttt{*.rrxiv.aux} file from which the rrxiv parser can produce
a CIR object that validates against \texttt{cir.schema.json}~\cite{rrxiv-cir-schema} v0.1.0.
\end{claim}

\begin{evidence}[CI conformance test]
\label{ev:ci-test}
The \texttt{tests/conformance/} suite in this repository runs the parser
on this file and asserts the result validates against the schema.
The test passes in CI on every push.
\end{evidence}

\begin{openquestion}[Future schema bumps]
\label{oq:schema-bumps}
The conformance contract above is anchored to \texttt{cir.schema.json}
v0.1.0. When the schema bumps, this fixture must continue to produce a
valid CIR \emph{or} be updated alongside the schema bump RRP.
\end{openquestion}

\bibliographystyle{plainnat}
\bibliography{minimal}

\end{document}
